\begin{table}[htbp]\centering\small
\caption{Marco macroeconómico, 2026--2027}\label{cua:macro}
\begin{tabular}{lrrr}
\toprule
Concepto & 2026 aprobado & 2026 estimado & \textbf{2027} \\
\midrule
PIB nominal (mmp) & 38.715,9 & 37.160,7 & \textbf{39.419,4} \\
Tipo de cambio promedio (pesos/dólar) & --- & 17,6 & \textbf{17,9} \\
Cetes 28 días, promedio (\%) & --- & 6,5 & \textbf{6,1} \\
Inflación INPC dic./dic. (\%) & 3,0 & 3,5 & \textbf{3,0} \\
Mezcla mexicana (dólares por barril) & 54,9 & 78,4 & \textbf{61,8} \\
Plataforma de producción (mbd) & 1.794,0 & 1.800,0 & \textbf{1.800,1} \\
Plataforma de exportación (mbd) & 521,0 & 522,4 & \textbf{426,6} \\
\bottomrule
\end{tabular}
\\[2pt]\begin{minipage}{0.92\textwidth}\scriptsize Fuente: CGPE 2027, edición SHCP, Anexos II.5 (p. 66) y III.1 (p. 68). El PIB de 2026 se revisa a la baja en 1.555,2 mmp entre el aprobado y el estimado; toda razón a PIB de 2026 depende de cuál añada se use.\end{minipage}
\end{table}
